\documentclass[11pt]{article}
\usepackage[margin=1in]{geometry}
\usepackage{amsmath}
\usepackage{amssymb}
\usepackage{tcolorbox}

\title{Problem 2b: Answer Format \\ (Exactly What to Write)}
\author{ECO 310 - The Format You Need}
\date{February 23, 2026}

\begin{document}

\maketitle

\section{What Problem 2b Asks}

\textbf{Question:} Write out the optimality conditions, \textit{if you guess that the utility constraint is binding, nonnegativity on good 1 is binding, and nonnegativity on good 2 is non-binding.}

\textbf{Translation:} Given this specific guess, what do the equations look like?

\section{The Guess in Math Terms}

The guess tells you:
\begin{itemize}
\item \textbf{Utility constraint is binding:} $x_1 + x_2 = u^*$ (exactly, not $>$)
\item \textbf{Nonnegativity on good 1 is binding:} $x_1 = 0$ (exactly)
\item \textbf{Nonnegativity on good 2 is non-binding:} $x_2 > 0$ (strictly positive)
\end{itemize}

\textbf{What this means for the $\lambda$ values:}
\begin{itemize}
\item Since $x_1 = 0$ (constraint binds), $\lambda_1 \geq 0$ (can be positive)
\item Since $x_2 > 0$ (constraint doesn't bind), $\lambda_2 = 0$ (complementary slackness!)
\item Since utility constraint binds, $\lambda_3 \geq 0$ (can be positive)
\end{itemize}

\section{Format of Your Answer}

Here's EXACTLY what you should write for Problem 2b:

\begin{tcolorbox}[colback=yellow!10!white,colframe=orange!75!black,title=Problem 2b Answer]

\textbf{Given the guess:}
\begin{itemize}
\item Utility constraint is binding: $x_1 + x_2 = u^*$
\item Nonnegativity on good 1 is binding: $x_1 = 0$
\item Nonnegativity on good 2 is non-binding: $x_2 > 0$
\end{itemize}

\textbf{The optimality conditions simplify to:}

\textbf{1. From $\frac{\partial \mathcal{L}}{\partial x_1}$ (since $x_1 = 0$):}
\begin{equation}
-p_1 + \lambda_1 + \lambda_3 \leq 0
\end{equation}

\textbf{2. From $\frac{\partial \mathcal{L}}{\partial x_2}$ (since $x_2 > 0$):}
\begin{equation}
-p_2 + \lambda_2 + \lambda_3 = 0
\end{equation}

Since $\lambda_2 = 0$ (complementary slackness with $x_2 > 0$):
\begin{equation}
-p_2 + \lambda_3 = 0 \quad \Rightarrow \quad \lambda_3 = p_2
\end{equation}

\textbf{3. From $\frac{\partial \mathcal{L}}{\partial \lambda_3}$ (utility constraint):}
\begin{equation}
x_1 + x_2 = u^*
\end{equation}

Since $x_1 = 0$:
\begin{equation}
x_2 = u^*
\end{equation}

\textbf{4. Solution from these conditions:}
\begin{align}
x_1 &= 0 \\
x_2 &= u^* \\
\lambda_1 &= p_1 - p_2 \quad \text{(from equation 1, substituting $\lambda_3 = p_2$)} \\
\lambda_2 &= 0 \\
\lambda_3 &= p_2
\end{align}

\textbf{5. For this to be valid, we need $\lambda_1 \geq 0$:}
\begin{equation}
\lambda_1 = p_1 - p_2 \geq 0 \quad \Rightarrow \quad p_1 \geq p_2
\end{equation}

This solution is valid when good 1 is at least as expensive as good 2.

\end{tcolorbox}

\section{Step-by-Step Logic}

Here's how you get to that answer:

\subsection{Step 1: Identify What You Know}

From the guess:
\begin{itemize}
\item $x_1 = 0$ (given)
\item $x_2 > 0$ (given)
\item $x_1 + x_2 = u^*$ (given)
\end{itemize}

From complementary slackness:
\begin{itemize}
\item Since $x_2 > 0$, we must have $\lambda_2 = 0$
\item Since $x_1 = 0$, we can have $\lambda_1 > 0$
\end{itemize}

\subsection{Step 2: Write the General Optimality Conditions}

From your Lagrangian $\mathcal{L} = -p_1 x_1 - p_2 x_2 + \lambda_1 x_1 + \lambda_2 x_2 + \lambda_3(x_1 + x_2 - u^*)$

General conditions are:
\begin{align}
\frac{\partial \mathcal{L}}{\partial x_1} &= -p_1 + \lambda_1 + \lambda_3 \leq 0, \quad x_1 \geq 0, \quad x_1(-p_1 + \lambda_1 + \lambda_3) = 0 \\
\frac{\partial \mathcal{L}}{\partial x_2} &= -p_2 + \lambda_2 + \lambda_3 \leq 0, \quad x_2 \geq 0, \quad x_2(-p_2 + \lambda_2 + \lambda_3) = 0 \\
\frac{\partial \mathcal{L}}{\partial \lambda_3} &= x_1 + x_2 - u^* \geq 0, \quad \lambda_3 \geq 0, \quad \lambda_3(x_1 + x_2 - u^*) = 0
\end{align}

\subsection{Step 3: Apply Your Guess to Simplify}

\textbf{For $x_1$ condition:}
\begin{itemize}
\item Since $x_1 = 0$, the complementary slackness $x_1(-p_1 + \lambda_1 + \lambda_3) = 0$ is satisfied automatically
\item We just need: $-p_1 + \lambda_1 + \lambda_3 \leq 0$
\end{itemize}

\textbf{For $x_2$ condition:}
\begin{itemize}
\item Since $x_2 > 0$, complementary slackness requires: $-p_2 + \lambda_2 + \lambda_3 = 0$ (exactly zero!)
\item Since we know $\lambda_2 = 0$: $-p_2 + \lambda_3 = 0$, so $\lambda_3 = p_2$
\end{itemize}

\textbf{For utility constraint:}
\begin{itemize}
\item Since constraint binds: $x_1 + x_2 = u^*$
\item Since $x_1 = 0$: $x_2 = u^*$
\end{itemize}

\subsection{Step 4: Solve for Everything}

Now substitute $\lambda_3 = p_2$ into the $x_1$ condition:
\begin{equation}
-p_1 + \lambda_1 + p_2 \leq 0 \quad \Rightarrow \quad \lambda_1 \leq p_1 - p_2
\end{equation}

For $\lambda_1 \geq 0$ (required), we need $p_1 - p_2 \geq 0$, or $p_1 \geq p_2$.

Actually, at the optimal solution, we typically have the inequality binding, so:
\begin{equation}
\lambda_1 = p_1 - p_2
\end{equation}

Wait, that gives $\lambda_1 \geq 0$ when $p_1 \geq p_2$. But this guess assumes we buy only good 2 ($x_1 = 0$, $x_2 = u^*$), which happens when good 2 is cheaper or equal price!

Let me recalculate...

If $\lambda_3 = p_2$ and $-p_1 + \lambda_1 + \lambda_3 \leq 0$:
\begin{equation}
\lambda_1 \leq p_1 - p_2
\end{equation}

Setting to equality (binding): $\lambda_1 = p_1 - p_2$

For this to be valid ($\lambda_1 \geq 0$): $p_1 \geq p_2$

Hmm, but intuitively if $x_1 = 0$ and $x_2 = u^*$, we're buying only good 2, which should happen when good 2 is cheaper ($p_2 < p_1$).

Actually, checking the inequality: $\lambda_1 = p_1 - p_2$. If $p_1 > p_2$, then $\lambda_1 > 0$, which is allowed.

\textbf{Conclusion:} The guess $x_1 = 0$, $x_2 = u^*$ is valid when $p_1 \geq p_2$ (good 1 is at least as expensive).

\section{What to Actually Write on Your Homework}

Just write what's in the yellow box above! That's the format your professor expects.

The key parts are:
\begin{enumerate}
\item State the guess clearly
\item Write simplified optimality conditions given that guess
\item Solve the system of equations
\item State the final solution: values of $x_1, x_2, \lambda_1, \lambda_2, \lambda_3$
\item Check validity (when does $\lambda_1 \geq 0$?)
\end{enumerate}

\section{Common Confusion}

\textbf{Q: Why does $x_2 > 0$ mean $\lambda_2 = 0$?}

\textbf{A:} Complementary slackness! The condition is $\lambda_2 \cdot x_2 = 0$. If $x_2 > 0$, then we must have $\lambda_2 = 0$ for the product to equal zero.

\textbf{Q: Why does $x_1 = 0$ mean $\lambda_1$ can be positive?}

\textbf{A:} Complementary slackness again: $\lambda_1 \cdot x_1 = 0$. If $x_1 = 0$, then $\lambda_1$ can be anything (including positive) and the product is still zero.

\textbf{Q: Where does $\lambda_3 = p_2$ come from?}

\textbf{A:} From the FOC for $x_2$: $-p_2 + \lambda_2 + \lambda_3 = 0$. Since $\lambda_2 = 0$, we get $\lambda_3 = p_2$.

\end{document}
